% ==============================================================================
% MODEL SECTION
% Welfare framework and mapping to the reduced form
% ==============================================================================

% ==============================================================================
% SLIDE 7
% ==============================================================================
\begin{frame}[t]
    \frametitle{Model setup: housing supply affects owners, renters, and potential entrants differently}
    \small

    \begin{columns}[T,totalwidth=\textwidth]
        \begin{column}{0.31\textwidth}
            \begin{block}{Owners}
                \footnotesize
                Local amenities, congestion, access, and asset values.
            \end{block}
        \end{column}
        \begin{column}{0.31\textwidth}
            \begin{block}{Renters}
                \footnotesize
                Rents, access/productivity, and local amenities.
            \end{block}
        \end{column}
        \begin{column}{0.31\textwidth}
            \begin{block}{Entrants}
                \footnotesize
                Gain when supply lets them enter New York rather than take the outside option.
            \end{block}
        \end{column}
    \end{columns}

    \vspace{0.4em}
    \begin{center}
        \begin{tikzpicture}[
            node distance=1.4cm,
            box/.style={draw=nycblue!60, fill=nycblue!6, rounded corners=2pt, text width=3.1cm, align=center, inner sep=6pt, font=\footnotesize},
            smallbox/.style={draw=nycgray!45, fill=nyclight, rounded corners=2pt, text width=2.5cm, align=center, inner sep=5pt, font=\scriptsize},
            arrow/.style={-{Stealth[length=2.1mm]}, thick, draw=nycgray}
        ]
            \node[box] (supply) {Approved housing\\\(q_j\)};
            \node[smallbox, above right=0.55cm and 1.2cm of supply] (rents) {Rents and\\location choice};
            \node[smallbox, right=1.4cm of supply] (wages) {Citywide\\agglomeration};
            \node[smallbox, below right=0.55cm and 1.2cm of supply] (costs) {Local\\development costs};
            \draw[arrow] (supply) -- (rents);
            \draw[arrow] (supply) -- (wages);
            \draw[arrow] (supply) -- (costs);
        \end{tikzpicture}
    \end{center}
\end{frame}

% ==============================================================================
% SLIDE 8
% ==============================================================================
\begin{frame}[t]
    \frametitle{Two regimes, three wedges}
    \small

    \begin{columns}[T,totalwidth=\textwidth]
        \begin{column}{0.48\textwidth}
            \begin{block}{Centralized approval}
                \begin{itemize}
                    \item Internalizes citywide and entrant benefits.
                    \item May mismeasure true neighborhood costs.
                \end{itemize}
            \end{block}
        \end{column}
        \begin{column}{0.48\textwidth}
            \begin{block}{Decentralized approval}
                \begin{itemize}
                    \item Uses local information about real costs.
                    \item May overweight incumbents and ignore spillovers.
                \end{itemize}
            \end{block}
        \end{column}
    \end{columns}

    \vspace{0.5em}
    \begin{center}
        \footnotesize
        \begin{tabular}{>{\raggedright\arraybackslash}p{0.25\textwidth}>{\raggedright\arraybackslash}p{0.62\textwidth}}
            \toprule
            \textbf{Wedge} & \textbf{Economic content} \\
            \midrule
            Externality & Benefits to other districts, potential entrants, citywide rents, labor markets, and agglomeration. \\
            Representation & Political weights differ from welfare weights; incumbent homeowners may be overweighted relative to renters and future residents. \\
            Information & Central authority may not observe true neighborhood costs. \\
            \bottomrule
        \end{tabular}
    \end{center}

\end{frame}

% ==============================================================================
% SLIDE 9
% ==============================================================================
\begin{frame}[t]
    \frametitle{Event-study estimates discipline the political approval wedge}
    \small

    \begin{columns}[T,totalwidth=\textwidth]
        \begin{column}{0.48\textwidth}
            \begin{block}{Approval wedge}
                \vspace{-0.8em}
                \[
                    p_{jt} = K'_j(q_{jt}) + \phi_{jt}
                \]
                \vspace{-1.0em}
                \[
                    \phi_{jt} =
                    \phi_j^0 + \lambda_t \rho T_j + u_{jt}
                \]
                \vspace{-1.1em}
                \scriptsize
                \begin{itemize}
                    \setlength\itemsep{0.15em}
                    \item \(\phi_{jt}\): political / regulatory approval wedge; \(T_j\): homeowner exposure.
                    \item \(\lambda_t\): local-control strength; \(\rho\): effect of exposure on the wedge.
                \end{itemize}
            \end{block}
        \end{column}
        \begin{column}{0.48\textwidth}
            \begin{block}{What the graph supplies}
                \scriptsize
                \begin{itemize}
                    \setlength\itemsep{0.25em}
                    \item \(\beta_t\): reduced-form homeowner gradient in observed supply.
                    \item 5+ unit buildings: negative gradient disciplines the sign and magnitude of \(\lambda_t\rho\).
                    \item Interpretation: local control binds more for discretionary projects that are likely to require rezoning or other special approvals
                \end{itemize}
            \end{block}
        \end{column}
    \end{columns}

    \vspace{0.15em}
    \begin{center}
        \begin{tikzpicture}[
            node distance=0.75cm,
            box/.style={draw=nycblue!60, fill=nycblue!7, rounded corners=2pt, text width=2.35cm, align=center, inner sep=3.5pt, font=\scriptsize\bfseries},
            arrow/.style={-{Stealth[length=2.1mm]}, thick, draw=nycgray}
        ]
            \node[box] (beta) {Homeowner\\gradient\\\(\beta_t\)};
            \node[box, right=of beta] (wedge) {Political\\approval wedge\\\(\lambda_t\rho\)};
            \node[box, right=of wedge] (cf) {Counterfactual land-use\\regimes};
            \draw[arrow] (beta) -- (wedge);
            \draw[arrow] (wedge) -- (cf);
        \end{tikzpicture}
    \end{center}
\end{frame}
